\documentclass{beamer}

\newcommand{\week}{Federal Reserve 2 of 5}

\title{History of the Federal Reserve}
\subtitle{Mishkin Chapter 13}
\author{Econ 357}
\date{\week}

% Reference the shared preamble
\input{preamble}

% Set the footer -- change 
\setbeamertemplate{footline}{
  \leavevmode%
  \vspace{2ex}
  \hbox{%
    % Left box: Econ 457
    \begin{beamercolorbox}[wd=.4\paperwidth,ht=2.5ex,dp=1ex,left]{author in head/foot}%
      \hspace{1em}Econ 357
    \end{beamercolorbox}%
    % Middle box: Week
    \begin{beamercolorbox}[wd=.2\paperwidth,ht=2.5ex,dp=1ex,center]{date in head/foot}%
      \centering\week
    \end{beamercolorbox}%
    % Right box: Slide numbers
    \begin{beamercolorbox}[wd=.4\paperwidth,ht=2.5ex,dp=1ex,center]{date in head/foot}%
      \hfill\insertframenumber{} 
    \end{beamercolorbox}%
  }%
  \vskip0pt%
}

\begin{document}

\frame{\titlepage}

\begin{frame}
    \frametitle{Federal Reserve}

    \textbf{Outline for Lectures on the Federal Reserve}
    \begin{enumerate}
        \item \fbox{History / Functions}
        \item Money Supply Process
        \item Dual Mandate
        \item Federal Open Market Committee (FOMC)
        \item Extraordinary Tools: Interest Rates and Balance Sheet
    \end{enumerate}
    \vspace{1em}
    Reference: Mishkin Chapters: 13, 14, 15, 16

\end{frame}

\begin{frame}
    \frametitle{Outline for Today's Lecture}
    
\begin{enumerate}
    \item Panic of 1907
    \begin{itemize}
        \item Financial Stability / Lender of Last Resort
    \end{itemize}
    \item Great Depression
    \begin{itemize}
        \item Monetary Policy
    \end{itemize}
    \item WWII
    \begin{itemize}
        \item  Financed US war effort
    \end{itemize}
    \item 1970s
    \begin{itemize}
        \item Volcker ends the inflationary period
    \end{itemize}
    
\end{enumerate}

\end{frame}

\begin{frame}[t]
    \frametitle{1. Panic of 1907}
    \framesubtitle{19th Century}

    In the early 1800s there were two separate iterations of a Central Bank.   Both times the bank was allowed to lapse, the second by Andrew Jackson in 1836.
    \vspace{1em}

    Absence of a central bank led to a chaotic financial system:
    \begin{itemize}
        \item In some periods there was too much money, which led to inflation and eventually a crash (e.g. 1840s)
        \item In other periods, there was not enough money, which led to deflation and depression (e.g. William Jennings Bryan crusade for Silver in the 1890s)
    \end{itemize}

    Even so, here was fierce resistance to reestablishing a central bank:
    \begin{itemize}
        \item Fear of centralized power
        \item Distrust of moneyed interests
    \end{itemize}

\end{frame}

\begin{frame}[t]
    \frametitle{1. Panic of 1907}
    \framesubtitle{JP Morgan}

    \begin{columns}
    
    \begin{column}{0.5\textwidth}
        \footnotesize
        \begin{itemize}
            \item Started with an attempted “short squeeze”, financed by borrowing.   The holders of the short positions withstood, and the resulting losses on the leveraged positions led to a bank run and eventual failure of NY’s 3rd largest bank.
            \item Contagion spread and interest rates between banks reached 70\%.
            \item Eventually stopped when JP Morgan organized a group financiers to support NY’s banks and stock exchange
        \end{itemize}
    \end{column}
    
    \begin{column}{0.5\textwidth}

        \centering
        \includegraphics[width=0.7\textwidth]{Econ357/figures/04_fed_1_jpm.png}
    \end{column}
    
    \end{columns}

\end{frame}

\begin{frame}[t]
    \frametitle{1. Panic of 1907}
    \framesubtitle{Financial Reform}

    \begin{itemize}
        \item The Panic of 1907 led bankers and legislators to prioritize financial reform, and eventually led to the Federal Reserve Act of 1913.
        \item The Fed was initially conceived to, in part, stabilize the financial system and stop banking panics.
        \item The Fed’s function was modeled after European central banks.   It was intended to be “lender of last resort” (the role that JP Morgan played in 1907).
        \item Bagehot’s dictum (1873): “Central banks should early and lend freely, against good collateral, and at high rates”
    \end{itemize}

\end{frame}

\begin{frame}[t]
    \frametitle{1. Panic of 1907}
    \framesubtitle{Creation of the Federal Reserve}

    \textbf{Federal Reserve Act of 1913:}
    \begin{itemize}
        \item Authority shared between the Board, located in DC, and twelve Regional Reserve Banks located throughout the country, mostly towards the East.
        \item Elaborate System of Checks and Balances, including that each Regional Bank was a quasi-public institution owned by private commercial banks in the district that are members of the Fed system
    \end{itemize}

\end{frame}

\begin{frame}[t]
    \frametitle{1. Panic of 1907}
    \framesubtitle{Creation of the Federal Reserve}

    \centering
    \includegraphics[width=0.8\textwidth]{Econ357/figures/04_fed_1_regions.png}

    \blfootnote{Source: Mishkin}
    
\end{frame}

\begin{frame}[t]
    \frametitle{2. The Great Depression}

    The Great Depression was the first Fed’s first big test with regards to monetary policy.   The Fed did not do well.
    \vspace{1em}

    \centering
    \includegraphics[width=0.7\textwidth]{Econ357/figures/04_fed_1_ue.png}

    \blfootnote{Source: Bureau of Labor Statistics}

\end{frame}

\begin{frame}[t]
    \frametitle{2. The Great Depression}

    \centering
    \includegraphics[width=0.6\textwidth]{Econ357/figures/04_fed_1_output.png}

    \blfootnote{Source: Friedman and Schwartz, “Monetary History of the United States”, 1963}

\end{frame}

\begin{frame}[t]
    \frametitle{2. The Great Depression}

    \centering
    \includegraphics[width=0.8\textwidth]{Econ357/figures/04_fed_1_ms.png}

    \blfootnote{Source: Friedman and Schwartz, “Monetary History of the United States”, 1963}

\end{frame}

\begin{frame}[t]
    \frametitle{2. The Great Depression}

    \begin{quote}
    "Regarding the Great Depression ... we did it. We're very sorry. ... We won't do it again."
    \end{quote}
    \vspace{1em}

    -- Ben Bernanke, November 8, 2002, in a speech given at "A Conference to Honor Milton Friedman ... On the Occasion of His 90th Birthday."

\end{frame}

\begin{frame}[t]
    \frametitle{3. World War II}

    During World War II the Federal Reserve actively bought US Treasury bonds in an effort to keep financing costs low and support the war.
    \vspace{1em}
    
    \centering
    \includegraphics[width=0.8\textwidth]{Econ357/figures/04_fed_1_bs.png}

    \blfootnote{Source: Chris Waller and Lowell Ricketts, “The Rise and (Eventual) Fall in the Fed's Balance Sheet”, 2014"}

\end{frame}

\begin{frame}[t]
    \frametitle{3. World War II}

    Inflation started to surge in 1950, due, in part, to the Korean War.   The situation was untenable for the Fed.
    \vspace{1em}
    
    \centering
    \includegraphics[width=0.8\textwidth]{Econ357/figures/04_fed_1_korea.png}



\end{frame}

\begin{frame}[t]
    \frametitle{3. World War II}

    \textbf{Fed-Treasury Accord of 1951:}
    \begin{itemize}
        \item The Fed ended its policy of capping government bond yields.
        \item The Fed-Treasury Accord removed the Treasury Secretary from the FOMC.  (Prior to 1951 the Treasury Secretary had been part of the Federal Open Market Committee)
        \item Started the period of Federal Reserve independence
    \end{itemize}

\end{frame}

\begin{frame}[t]
    \frametitle{4. 1970s}
    \framesubtitle{Inflation}

    There were a number of inflationary pressures in the late 1960s and 1970s:
    \begin{itemize}
        \item Large budget deficits due to Vietnam War and Great Society programs
        \item Oil price rose sharply due to: OPEC Embargo (1973) and Fall of the Shah of Iran (1979)
        \item Money supply growth was steady and high
    \end{itemize}

\end{frame}

\begin{frame}[t]
    \frametitle{4. 1970s}
    \framesubtitle{Political Influence on the Fed}

    During this time, both President’s Johnson and Nixon pressured the Fed to maintain easy policy.   This likely contributed to the inflation problem.
    \begin{itemize}
        \item \textbf{Johnson}: Called the Fed chair to his ranch in Texas in 1965, “shoved him against the wall” and told him to cut rates.
        \item \textbf{Nixon}: Told the Fed chair that policy needed to be easier to help him be re-elected in 1972 (was recorded in the Oval Office saying that)   Put public pressure the Fed Chair by releasing his salary.
    \end{itemize}

\end{frame}

\begin{frame}[t]
    \frametitle{4. 1970s}
    \framesubtitle{Inflation}

    \centering
    \includegraphics[width=\textwidth]{Econ357/figures/04_fed_1_infl1970s.png}

\end{frame}

\begin{frame}[t]
    \frametitle{4. 1970s}
    \framesubtitle{Paul Volcker}

    \begin{columns}
        \begin{column}{0.5\textwidth}
            \begin{itemize}
                \item Appointed by President Carter in 1979.   Reappointed by President Reagan.
                \item Credited with ending the inflationary period of the 1970s.
                \item Reappointed by Reagan, even though the US economy was in a recession in 1982
            \end{itemize}
        \end{column}
        \begin{column}{0.5\textwidth}
            \centering
            \includegraphics[width=0.6\textwidth]{Econ357/figures/04_fed_1_volcker.png}
        \end{column}
    \end{columns}

\end{frame}

\begin{frame}[t]
    \frametitle{4. 1970s}
    \framesubtitle{Fed Funds}

    \centering
    \includegraphics[width=\textwidth]{Econ357/figures/04_fed_1_fedfunds.png}

\end{frame}

\begin{frame}[t]
    \frametitle{4. 1970s}
    \framesubtitle{Unemployment Rate}

    \centering
    \includegraphics[width=\textwidth]{Econ357/figures/04_fed_1_ue1970s.png}

\end{frame}

\begin{frame}[t]
    \frametitle{Summary}

    \begin{itemize}
        \item The Fed was created in 1913, as a response to the Panic of 1907.   The Fed was initially created to stabilize the financial system and to be a “lender of last resort”
        \item The Fed contributed to the severity of the Great Depression by allowing the money supply to contract.
        \item The Fed supported the government in WWII by buying US Treasury bonds.   While effective during the war, this policy contributed to post-war inflation.
        \item The Fed badly mis-managed inflation in the 1970s.   Chairman Paul Volcker ended the inflationary period by causing a series of severe recessions in the early 1980s.
    \end{itemize}

\end{frame}

\end{document}